% Auto-extracted table snippet; requires \usepackage{booktabs}.
% Caption: TABLEIII PERFORMANCE EVALUATIONOF MEAN OPTION SCORE (USER STUDY) ACROSS VARIOUS VIDEO ENHANCEMENT METHODS
\begin{tabular}{lll}
\toprule
 & Sequences & STDF-R3 QECF CAT TGAF LSFMD \\
\midrule
E & ChineseEditing scwebbrowsing nglishDocumentEditing scmap scprogramming & 50.981 40.364 39.907 48.659 61.938 44.173 54.943 50.282 58.994 64.764 39.741 37.361 54.422 45.547 59.476 52.067 52.233 56.432 51.794 57.057 54.529 48.091 51.833 45.942 57.343 \\
 & itches: To evaluate the capability of our proposed LSFMD handling dramatic motion and scene switches, two different es of screen content videos were selected to compute ∆PSNR curves for STDF-R3, QECF, CAT, TGAF, STA, -STIF-M, STDR, and our proposed method. The scpro- mming sequence involves pop-up windows and window itching. These dynamic motions are commonly seen in daily and can pose difficulties for video quality enhancement al- ithms. Additionally, the scSlideShow sequence is composed spliced videos from CTC [34], allowing us to evaluate the formance of our method in scenarios involving abrupt scene nsitions. The results are shown in Fig. 9, where dashed es indicate scene switch frames and gray shadow regions tinguish the frames exhibiting dynamic motion. The result Fig. 9(a) demonstrates that our proposed LSFMD mostly performs the others from frame 38 to frame 63 in the rogramming sequence. This shadow region encompasses ndow switches and a pop-up window. It can demonstrate t our proposed method can achieve significant ∆PSNR ing periods of dramatic motion. While the STA only utilizes single frame to handle the scene switch which does not form well in dramatic motion. In Fig. 9(b), frame 28 and me 44 represent the switch points between two PowerPoint des in the scSlideShow sequence. Notably, our proposed thod demonstrates an improvement during most of the nsition points, highlighting its effectiveness in handling upt scene transitions. In summary, our approach can take balance between the performance in dynamic content and &  \\
\bottomrule
\end{tabular}
